% Körper – bank/koerper/ (zusammenbau v0.4, Heft)
\einheitenkopf[t2]{Körper}
\setcounter{aufgabe}{1}

\begin{aufgabe}{Körper und Netze – Prüfungsaufgaben}
\begin{teile}
\teil Ein Zelt hat diese Form. Kreuze an, welcher Körper es ist \hfill (P10 2017 OS). \\ \kreuz{Pyramide} \\ \kreuz{Prisma} \\ \kreuz{Quader}

\prismadreieck{4}{2}{2.5}{}{}{}
\teil Ein Stück Käse hat diese Form. Kreuze an, welcher Körper es ist \hfill (P10 2017 OS). \\ \kreuz{Quader} \\ \kreuz{Pyramide} \\ \kreuz{Prisma}

\prismadreieck{3}{3}{1.5}{}{}{}
\teil Aus diesem Netz wird ein Körper gefaltet. Kreuze an, welcher Körper entsteht \hfill (P10 2018 OS). \\ \kreuz{Prisma} \\ \kreuz{Pyramide} \\ \kreuz{Quader}

\netzpyramide{2.5}{2}{}{}
\end{teile}
\end{aufgabe}

\begin{aufgabe}{Körper und Netze – Prüfungsaufgaben – weiter}
\begin{teile}
\teil Aus dem Netz wird eine Geschenkschachtel in Würfelform gefaltet. Die Fläche D ist der Deckel. Gib an, welche Fläche der Boden ist \hfill (P10 2016 OS).

$\begin{array}{cccc}  &  & A &  \\ D & N & M &  \\  &  & B & C \end{array}$
\leerfeld
\teil Das Netz wird zu einem Würfel gefaltet, der Würfel wird auf die Fläche D gestellt. Gib an, welche Fläche dann oben liegt \hfill (P10 2016 OS).

$\begin{array}{cccc}  & A &  &  \\  & M & N & D \\ C & B &  &  \end{array}$
\leerfeld
\teil Ein Pavillondach hat die Form einer Pyramide mit sechseckiger Grundfläche. Gib an, wie viele Kanten diese Pyramide hat \hfill (P10 2019 OS). \leerfeld Kanten
\teil Ein Briefbeschwerer hat die Form einer Pyramide mit dreieckiger Grundfläche. Gib an, wie viele Kanten diese Pyramide hat \hfill (P10 2019 OS). \leerfeld Kanten
\end{teile}
\end{aufgabe}

\begin{aufgabe}{Quader – Prüfungsaufgaben}
\begin{teile}
\teil Ein Würfel hat die Kantenlänge $8$ cm. Gib sein Volumen an \hfill (P10 2026 FOR). \leerfeld[cm³]
\teil Ein würfelförmiger Blumenkübel ist innen $2$ dm lang, breit und hoch. Gib an, wie viel Erde hineinpasst \hfill (P10 2026 FOR). \leerfeld dm³
\end{teile}
\end{aufgabe}
